\paragraph{Crossing and completion ranges.} Prediction 1 held: $\Delta(e)$
changes sign between .93 and .94 and its interpolated root is .93361, with
bootstrap interval [.93358, .93365]. All 20,000 draws had exactly one crossing
on the evaluated grid. This supersedes the earlier two-point interpolation
between $e=.9$ and 1 as a numerical location in the original workflow, but the
bootstrap interval is not a bound on the true root: interpolation error is not
included, and it is larger than the interval. $\Delta(e)$ is concave near the
root (second differences $-0.30$ and $-0.33$ points at .93 and .94 on the
0.01 grid), so the linear root lies to the left of the smooth one; quadratic,
cubic and quartic fits to neighbouring grid values give .93379 to .93382
(post hoc), outside the bootstrap interval. On the isolated bridge, where the
root is known exactly, the same linear method on the same grid gives .92331
against .92353, a miss of $2.2\times10^{-4}$, about three times the interval's
width; a quadratic through neighbouring grid values comes within
$2\times10^{-5}$. The crossing is therefore best stated as about .934 (smooth
fits .9338); Prediction 1 holds under either estimate. $\Delta(e)$ is strongly
nonlinear in $e$: $+3.9$ points at $e=.5$,
largest ($+9.0$) at $e=.8$, and $\mn21.5$ at $e=1$
(Table~\ref{tab:extension-crossing}), which is why the two-point interpolation
understated the root. The fine-grid crossing lies above the isolated
two-opportunity root .92353 and near the renewal comparator's .936. Per-length
isolated roots (.845 for four opportunities) do not bound it, because the two
arms weight bridge lengths differently: 26.2\% of the clustered baseline
coefficient, but 1.0\% of the independent one, comes from bridges with four or
more opportunities.

